\chapter*{\centering ABSTRAK}
\addcontentsline{toc}{chapter}{ABSTRAK}

\begin{center}
\setstretch{1.2}
\begin{tabular}{l l}
Nama          & : Gilang Pratama Putra Siswanto \\ 
Program Studi & : Fisika \\ 
Judul         & : \begin{tabular}[c]{@{}l@{}}Perancangan Kit Eksperimen Otomatisasi untuk \\ Hukum Pergeseran Wien Berbasis Sensor RGB \\ TCS34725, Arduino Uno R3, dan Antarmuka Python 3\end{tabular} \\ 
\end{tabular}
\end{center}

\vspace{1cm}

\begin{spacing}{1.25}
\noindent Penelitian ini bertujuan untuk menganalisis akurasi dan ketelitian kit eksperimen Hukum Pergeseran Wien berbasis Arduino Uno dan sensor TCS34725 dalam mengukur hubungan antara suhu dan panjang gelombang radiasi benda hitam. Dengan mengadopsi teknologi modern, eksperimen ini dilakukan di Laboratorium Instrumentasi-Robotika Universitas Islam Negeri Sunan Gunung Djati pada bulan September-Oktober 2024. Metodologi penelitian meliputi penggunaan sensor suhu dan sensor spektrum yang terintegrasi dengan mikrokontroler untuk menentukan panjang gelombang maksimum yang dipancarkan oleh benda pada suhu tertentu. Hasil pengujian menunjukkan bahwa sistem digital yang dikembangkan dapat bekerja secara optimal dan akurat, menghasilkan nilai konstanta pergeseran Wien (b) sebesar 0,002897773 dengan ketelitian 1,8796E-09 dan ketepatan 99,99\% terhadap nilai literatur. Penelitian ini diharapkan memberikan pemahaman lebih mendalam tentang aplikasi otomatisasi dalam eksperimen fisika serta kontribusinya terhadap pengembangan ilmu pengetahuan. \\

\noindent\textbf{Kata kunci:} Hukum Pergeseran Wien, Arduino Uno, sensor TCS34725, pengukuran suhu, radiasi benda hitam.
\end{spacing}
